\section*{Introduzione}

Appunti ordinati, con approfondimenti passo-passo, del corso di Circuiti ed algoritmi per il Digital Signal Processing per il corso di laurea in Ingegneria Elettronica 
presso l’Università Politecnica delle Marche per l'A.A. 2025/2026. \newline
        
Per quelli che seguono il nuovo corso di Ingegneria elettronica: \\
\url{https://www.univpm.it/Entra/Offerta_formativa_1/Corso_di_laurea_triennale_in_Ingegneria_Elettronica_e_delle_Tecnologie_Digitali} \\
cioè la Laurea in Ingegneria Elettronica e delle Tecnologie Digitali, 
il corso prende il nome di Digital Signal Processing. \newline 

Il corso di DSP è stato molto utile perchè mi ha dato una preparazione "matematica" su come i calcolatori (cioè computer e micro-controllori per la maggior parte) capiscono e interpretano il mondo reale e su come "simulando" i circuiti nel mondo digitale, 
si possono ottimizzare le elaborazioni per il mondo reale (questa frase ha qualcosa di filosofico, ma sarò un futuro ingegnere, quindi non ho le competenze per parlare di filosofia). \newline 

I principi matematici sarebbero tanti, gli approfondimenti pure, ma, a causa dei pochi CFU che questo esame ha, durante le lezioni non si apprfondiscono le tematiche e, purtroppo, bisogna fare un principio di fede agli argomenti. \newline 

Molti degli argomenti, però, verranno approfonditi alla magistrale; quindi, per chi fosse realmente interessato alle tematiche, vi consiglio di approfondire i prossimi anni. \newline 

Oppure, se volete già approfondire le tematiche adesso, potete consultare il libro "Discrete-Time Signal Processing" di Alan Oppenheim e Ronald Schafer (qualsiasi edizione va più che bene). \newline 

Un libro fantastico dove non si lascia nulla al caso. \newline 

Personalmente, per questo esame, ve lo sconsiglio perchè l'anno scorso ho provato a seguire sia le lezioni che il libro, ma una lezione del professore erano, almeno, 50 pagine del libro. \newline 

Se siete realmenti amanti di questa materia, studiate prima dalle dispense e dagli appunti presi a lezione (se avete la possibilità di venirci perchè, come ben so, alcuni di voi che mi stanno leggendo già lavorano); 
se non avete capito bene il tema, approfondite prima dal libro e poi chiedete ricevimento. \newline 

Come scritto anche in altri appunti, e se mi conoscete anche personalmente sapete che è una cosa che ripeto sempre, chiedete ricevimento al professore, almeno uno prima di un esame: 
meglio fare "lo stupido" durante il ricevimento, che non riuscire a scrivere la risposta alla domanda in sede di esame. \newline 

I professori del nostro corso sono, previo avviso via email, sempre disponibili per un ricevimento ed è una risorsa preziosa che, molti di noi, non sfruttano abbastanza e che abbiamo in abbondanza. \newline 

Andateci, non siate timidi, e andateci preparando prima bene le domande (meglio domande secche se possibile). \newline 

Ogni giorno, mi confronto con colleghi di altri corsi e di altri atenei, ed il ricevimento è una risorsa molto rara, che in molti vorrebbero, ma che non possono avere perchè i loro professori non possono (magari perchè hanno tanti incarichi nella vita accademica o anche nella vita privata o tutti e due). \newline 

In generale, se siete in difficoltà, chiedete una mano. \newline 

Ritornando al DSP, esso ha diversi campi di utilizzo: dall'ambito musicale (quello che interessa maggiormente al professore), a quello delle telecomunicazioni (come ad esempio, vedendo su LinkedIn i miei futuri colleghi, il DSP trova tanto ambito, visto anche il periodo storico in cui ci troviamo adesso, 
nei radar, sia civili che militari). \newline 

Per quanto riguarda questi appunti, ovviamente, non tutto è "farina del mio sacco": le fonti sono sempre indicate e voglio dare credito a loro. \newline 

Non voglio essere un sostituto al professore, ma un compagno, come voi, che ha studiato, si è rimboccato le maniche e che racconta un po' quello che gli piace fare di più, 
cioè studiare ingegneria per poi applicare i concetti teorici nella realtà. \newline 

Inoltre, siccome questi appunti si trovano nel mio GitHub al seguente link \\
\href{https://github.com/ciccio25/appunti-dsp}{https://github.com/ciccio25/appunti-dsp} \\
e pubblici, se trovate qualche errore o volete cotribuire per migliorarlo, perchè banalmente ci sono degli "orrori grammaticali" o ci sono delle imprecisioni, 
mi potete contattare telefonicamente (se avete il mio numero di cellulare), o aprire in "Issue" qui su GitHub o mandarmi una email a \\
\href{mailto:rossini.stefano.appunti@gmail.com}{rossini.stefano.appunti@gmail.com}. \newline 

Ultima cosa: se proprio volete stampare questi appunti, vi consiglio di stamparli a colori, perchè i colori differenti, in diverse immagini, aiutano molto nella compresione della figura stessa. \newline 

Basta con le chiacchiere. \newline 

Buono studio e buona vita. \newline 

\newpage 

